%%%%%%%%%%%%%%%%%%%%%%%%%%%%%%%%%%%%

\section{ケーススタディ}

%%%%%%%%%%%%%%%%%%%%%%%%%%%%%%%%%%%%

\begin{frame}
\frametitle{慢性疲労症候群の治療}

\begin{itemize}

\item \textbf{目的：}慢性疲労症候群に対する認知行動療法の効果を評価する。

\item \textbf{参加者候補：}かかりつけ医や専門医から慢性疲労症候群専門クリニックへ紹介された142名の患者。

\item \textbf{実際の参加者：}紹介された142名のうち、実際に研究へ参加したのは60名のみ。
診断基準を満たさない者、他の健康上の問題を抱える者、参加を拒否した者は除外された。

\end{itemize}

\ct{Deale et. al. \textit{Cognitive behavior therapy for chronic fatigue syndrome: A randomized controlled trial}. The American Journal of Psychiatry 154.3 (1997).}

\end{frame}

%%%%%%%%%%%%%%%%%%%%%%%%%%%%%%%%%%%%

\begin{frame}
\frametitle{研究デザイン}

\begin{itemize}

\item 患者は治療群と対照群にランダムに割り当てられ、各群30名：
\begin{itemize}
\item \hl{治療群}：認知行動療法――協調的・教育的なアプローチで行動面を重視。
症状を悪化させることなく活動量を段階的・安全に増やす方法を患者に指導した。
\item \hl{対照群}：リラクゼーション――活動量を増やす指導は行わず、
筋弛緩法・視覚化・即時リラクゼーション技法を指導した。
\end{itemize}

\end{itemize}

\end{frame}

%%%%%%%%%%%%%%%%%%%%%%%%%%%%%%%%%%%%

\begin{frame}
\frametitle{結果}

下の表は、6か月後の追跡調査における良好な転帰を示した患者の分布を示す。
なお、7名（治療群3名・対照群4名）が途中で脱落した。

\begin{center}
\begin{tabular}{ll  cc c}
			&				& \multicolumn{2}{c}{\textit{良好な転帰}} \\
\cline{3-4}
			&							& あり 	& なし 	& 合計	\\
\cline{2-5}
						&治療群 	& 19	 	& 8		& 27 	\\
\raisebox{1.5ex}[0pt]{\textit{群}}	&対照群		& 5	 	& 21	 	& 26 \\
\cline{2-5}
						&合計		& 24		& 29		& 53
\end{tabular}
\end{center}

\begin{itemize}

\item 治療群における良好な転帰の割合：
\[ 19 / 27 \approx 0.70 \rightarrow 70\% \]

\item 対照群における良好な転帰の割合：
\[ 5 / 26 \approx 0.19 \rightarrow 19\% \]

\end{itemize}

\end{frame}

%%%%%%%%%%%%%%%%%%%%%%%%%%%%%%%%%%%%

\begin{frame}
\frametitle{結果の解釈}

\dq{データはグループ間の「真の」差を示しているか？}

\begin{itemize}

\item コインを100回投げることを考えてみよう。
表が出る確率は50\%だが、ちょうど50回表が出るとは限らない。
このようなばらつきは、ほぼあらゆるデータ生成プロセスで生じる。

\item 2グループ間で観測された差（70 - 19 = 51\%）は真の差かもしれないし、
自然なばらつきによるものかもしれない。

\item 差が非常に大きいため、真の差である可能性が高いと考えられる。

\item この差が偶然によるものでないと判断できるかどうかは、統計的手法を用いて判断する必要がある。

\end{itemize}

\end{frame}

%%%%%%%%%%%%%%%%%%%%%%%%%%%%%%%%%%%%

\begin{frame}
\frametitle{結果の一般化}

\dq{この研究の結果は、慢性疲労症候群の患者全体に一般化できるか？}

これらの患者は特定の特性を持ち、自ら参加を希望した者たちであるため、
慢性疲労症候群の患者全体を代表しているとは言い難い。
ただちに全患者に一般化することはできないが、
この最初の研究は有望な結果を示している。
特定の特性を持つ患者には効果があることが示されており、
他の患者にも少なからず効果が期待できる。

\end{frame}

%%%%%%%%%%%%%%%%%%%%%%%%%%%%%%%%%%%%
